\documentclass[conference]{IEEEtran}
\IEEEoverridecommandlockouts
% The preceding line is only needed to identify funding in the first footnote. If that is unneeded, please comment it out.
\usepackage{cite}
\usepackage{amsmath,amssymb,amsfonts}
\usepackage{algorithmic}
\usepackage{graphicx}
\usepackage{textcomp}
\usepackage{xcolor}
\usepackage{tabularx}
\usepackage{array}
\usepackage{etoolbox}
\usepackage{float}
\usepackage{flafter}
\usepackage{placeins}
\usepackage{stfloats}        % lets figure* use page bottoms too
\raggedbottom                % stop columns stretching to fill
\setlength{\textfloatsep}{10pt plus 2pt minus 2pt}
\setlength{\floatsep}{8pt plus 2pt minus 2pt}

% Display math one step smaller, with tighter skips above and below.
\AtBeginEnvironment{equation}{\small}
\makeatletter
\g@addto@macro\normalsize{%
  \setlength\abovedisplayskip{5pt plus 2pt minus 2pt}%
  \setlength\belowdisplayskip{5pt plus 2pt minus 2pt}%
  \setlength\abovedisplayshortskip{3pt plus 1pt}%
  \setlength\belowdisplayshortskip{3pt plus 1pt}%
}
\makeatother

\def\BibTeX{{\rm B\kern-.05em{\sc i\kern-.025em b}\kern-.08em
    T\kern-.1667em\lower.7ex\hbox{E}\kern-.125emX}}
\begin{document}

\title{Characterizing the Robustness of Musical Instrument Classification Under Acoustic Perturbations}


\author{
\IEEEauthorblockN{1 \textsuperscript{st} Allan Luo Yu}
\IEEEauthorblockA{\textit{Fairview High School} \\
Boulder, USA \\
0009-0001-1403-2316}
\and
\IEEEauthorblockN{1\textsuperscript{st} Max Liu}
\IEEEauthorblockA{\textit{Clements High School} \\
Sugar Land, USA \\
MaxLiu2k@gmail.com}
\and
\IEEEauthorblockN{2\textsuperscript{nd} Gavin Hu}
\IEEEauthorblockA{\textit{St. Marks School} \\
Southborough, USA \\
gavinhu42@gmail.com}
}

\maketitle

\begin{abstract}
Instrument detection underpins Music Information Retrieval (MIR) infrastructure and supports applications including musicology, audio editing, and transcription. However, instrument recognition deep learning models may fail or lose substantial accuracy under artificial and real-world noise. This work benchmarks six instrument classification models---three trained from scratch (SVM, CNN, CRNN) and three pretrained (AST, MERT, and PANNs)---under varying forms and degrees of noise. We use the Philharmonia Orchestra Sound Samples dataset, trimmed to 12 common orchestral instruments. Noise was white (Gaussian), human non-speech from ESC-50, and environmental from DEMAND, at SNRs from 50 dB to $-$10 dB. All six models were scored on byte-identical mixtures, enabling paired comparison.

We find that our best clean (without noise) macro-F1 baselines are AST (0.991) and PANNs (0.989), followed by MERT (0.980), SVM (0.979), CRNN (0.974), and CNN (0.971). On clean audio all six models fell within two points of one another, but robustness diverged sharply: at 20 dB SNR the SVM retained only 10.4\% of its clean macro-F1 under white noise against 68.0\% under human non-speech, and model ranking inverted—SVM fell from fourth on clean audio to last under white noise. Clean macro-F1 is therefore not indicative of noise robustness, and degradation depends on noise category, not severity alone.
\end{abstract}

\begin{IEEEkeywords}
Audio Classification, Environmental Noise, Robustness, Music Information Retrieval, SNR, Deep Learning
\end{IEEEkeywords}

\section{Introduction}
Musical instruments are integral to the production and characterization of music, and their distinct timbral properties are intertwined with music composition, performance, and perception [1]. Consequently, automatically classifying instruments has become an active area of research, with applications in music transcription, source separation, post-production and automatic mixing, and music recommendation [2].

Early instrument classification systems ($\sim$1999) relied on statistical classifiers such as Gaussian Mixture Models (GMM) and Support Vector Machines (SVM) [3]. In the early 2000s, performance was improved with high-quality feature representation [4], and also drew attention to the difficulties of real-world recordings [5]. From around 2015, convolutional neural networks (CNNs) transformed instrument classification, and subsequent works pushed accuracy and precision past 95\% on clean, isolated-note datasets [6, 7, 8].

High accuracy on clean instrument audio reveals little about performance under perturbation; controlled degradation probes that relationship directly. Other instrument classification research shifts towards realism, such as polyphonic recognition [9], whereas noise-robust classification is much less explored. Where robustness analysis has been done, model span is narrow. Deng et al. [10] studied the pretrained models OpenL3, PANNs, and CLAP under parameterized audio effects (gain, low-pass filtering, reverberation, and bitcrushing), finding that pretraining does not automatically confer robustness; additive noise, the perturbation class studied here, was not covered. Elsewhere, comparisons similarly span only a single model family, such as pretrained embeddings [11]. No prior instrument classification study was identified that spans handcrafted features, models trained from scratch, and pretrained backbones. This breadth is deliberate but carries a trade-off: pretraining corpora, input representation, sample rate, architecture, and whether the backbone is fine-tuned or frozen all differ, so family-level differences cannot be attributed to a single factor. Complete classification systems are therefore compared, and performance differences are not attributed to architecture alone.

The research question is: how do instrument classification systems differ in their ability to retain classification performance under white, human, and environmental additive noise across a range of signal-to-noise ratios? The contributions are: (i) a matched robustness benchmark of six model families spanning handcrafted, scratch-trained, and pretrained approaches; (ii) an evaluation protocol in which every model is scored on the same generated noisy mixtures, enabling paired comparison; and (iii) per-instrument and per-noise degradation analysis, including a pre-registered test of whether acoustically similar instruments are confused more heavily under noise.


\section{Methodology}

\subsection{Materials}
The Philharmonia Orchestra's sound samples dataset was used [12], which comprises thousands of samples across over 20 instruments. Specifically, 12 common orchestral instruments were analyzed—bassoon, cello, clarinet, double bass, flute, french horn, oboe, trombone, trumpet, tuba, viola, and violin. Instruments averaged 45.3 distinct pitches. Class imbalance ranges from 433 (trumpet) to 852 (violin), a ratio of 1.97:1. Articulation was filtered to prevent further class imbalance. Ultimately, the dataset had 8,374 unique recordings. For human non-speech and environmental noise, the ESC-50 [13] and DEMAND [14] datasets were used.

\subsection{Data Pre-Processing \& Splitting}
\subsubsection{Pre-Processing}
Leading and trailing silence more than 30 dB below each recording’s peak was trimmed. One 3.0-second window was created per source recording, tiling (repeating) all recordings under 3.0 seconds. Audio was resampled to 22.05 kHz. Loudness was normalized by scaling each window towards an RMS amplitude of 0.1, with gain capped to keep amplitude below 0.99.
\subsubsection{Splitting}
Recordings were split by instrument-pitch group, targeting 70/15/15 train, validation, and test, so that performances of the same note at different dynamics could not appear in more than one split. The split was deterministic, yielding 5,861 training, 1,258 validation, and 1,255 test windows across 544 pitch groups, and was reused by all models.

\subsection{Modeling}

\begin{table*}[t]
\centering
\caption{Model configurations. B: batch size; E: maximum epochs;
ES: early-stopping patience; CE: cross-entropy.}
\label{tab:model_configs}
\footnotesize
\setlength{\tabcolsep}{5pt}
\begin{tabularx}{\textwidth}{
    @{}
    >{\centering\arraybackslash}p{1.4cm}
    >{\centering\arraybackslash}p{3.7cm}
    >{\centering\arraybackslash}X
    >{\centering\arraybackslash}p{3.6cm}
    @{}
}
\hline
\textbf{Model} &
\textbf{Input} &
\textbf{Training} &
\textbf{Schedule} \\
\hline

SVM &
88-D handcrafted features &
RBF; $C=100$, $\gamma=0.001$ &
Fixed run \\

CNN &
128-band log-Mel, 22.05 kHz &
Scratch; weighted CE; AdamW, $10^{-3}$ &
B32; E40; ES8; 5 seeds \\

CRNN &
128-band log-Mel, 22.05 kHz &
Scratch; weighted CE; AdamW, $10^{-3}$ &
B32; E40; ES8; 5 seeds \\

MERT &
Waveform, 24 kHz &
MERT-v1-95M fine-tuned; CE; AdamW,
$3{\times}10^{-5}$ backbone and $5{\times}10^{-3}$ head &
B16; E45; ES10; 1 seed \\

AST &
Log-Mel, 16 kHz &
AudioSet pretrained; weighted CE; AdamW, $10^{-5}$ &
B8; E10; 1 seed \\

PANNs &
Waveform, 32 kHz &
AudioSet CNN14 fine-tuned; weighted CE; AdamW, $10^{-4}$ &
B16; E25; ES10; 1 seed \\
\hline
\end{tabularx}
\end{table*}

Six models spanning a range of learning methods are compared. Because architectures differ substantially, each was trained with its native recipe (Table~\ref{tab:model_configs}). Weights were fitted on the training split; hyperparameters and checkpoints were selected on validation only.

\textbf{SVM.} Frame-wise acoustic features were summarized by their
temporal means and standard deviations into 88 dimensions per window,
normalized using training-set statistics.
\textbf{MERT.} MERT-v1-95M [15] returns 13 hidden-state
representations, each time-averaged, combined by learned layer weights
and passed to a 12-class linear head. The backbone and head were
fine-tuned jointly.
\textbf{AST.} Initialized from
MIT/ast-finetuned-audioset-10-10-0.4593 [16]; Log-Mel spectrograms,
backbone and head fine-tuned jointly.
\textbf{PANNs.} AudioSet-pretrained CNN14 [17] computes log-Mel
internally. The backbone was fine-tuned and its classifier replaced by
a 12-way head over a 2048-dimensional embedding.
\textbf{CNN.} Three convolutional blocks of 32, 64, and 128 channels
($3\times3$ convolution, batch normalization, ReLU, $2\times2$ max
pooling; $\sim$111{,}000 parameters).
\textbf{CRNN.} The same blocks as CNN, except the final one pools only on
frequency. Frequency is then collapsed to a 32-step sequence of
128-dimensional vectors read by a bidirectional GRU (128 units per
direction), time-averaged and passed to a linear classifier
($\sim$294{,}000 parameters).

Neural model training used validation-based checkpoint selection, early stopping, dropout, weight decay, and pretrained initialization where applicable. CNN and CRNN predictions were ensembled across 5 seeds (42-46).


\subsection{Noise Construction}
All models were evaluated on the same mixtures. Three noise categories were used: white (Gaussian, spectrally flat), human non-speech, and environmental. Human non-speech noises were drawn from ESC-50 targets 20–29. Environmental noise was drawn from all 18 DEMAND environments, spanning domestic, nature, office, public, street, and transportation settings. For each ESC-50 and DEMAND realization, one source clip or environment was drawn at random and one random three-second segment extracted, rejecting near-silent extracts, then resampled to 22.05 kHz.

Let $x$ denote a clean waveform and $n$ a selected noise waveform, each containing $N$ samples. The noise waveform was first centered by subtracting its mean:
\begin{equation}
\bar{n} = \frac{1}{N}\sum_{i=1}^{N}n_i,
\qquad
n_{c,i} = n_i-\bar{n},
\label{eq:noise_centering}
\end{equation}
where $n_c$ is the centered noise waveform. The power $P(z)$ of any waveform $z$ was defined as its mean squared amplitude. For a target signal-to-noise ratio $s$ in decibels, the centered noise was multiplied by gain
\begin{equation}
\alpha =
\sqrt{
\frac{P(x)}
{P(n_c)\,10^{s/10}}
}.
\label{eq:noise_gain}
\end{equation}
The noisy mixture was then constructed as
\begin{equation}
y_i = x_i+\alpha n_{c,i},
\label{eq:noise_mixture}
\end{equation}
where $y$ is the final noisy waveform.

Mixtures were generated at eight SNRs (50, 40, 30, 20, 10, 0, $-5$, and $-10$ dB). Two independently seeded noise draws were generated for each of the 1,255 test recordings and noise categories, producing 60,240 noisy WAV files in total ($1{,}255 \times 3\,\text{ noise categories} \times 8\,\text{ SNRs} \times 2\,\text{ realizations}$). Within each replicate, a single realization was reused and rescaled across all eight SNRs, so change along a robustness curve could be attributed to noise level instead of changing noise.

Each realization was seeded deterministically from the dataset identity, window, noise category and replicate index, making the corpus reproducible. After writing, each mixture was reloaded and required to match the target SNR within 0.1 dB, with residual DC below 1\% of the added-noise power. Mixtures were stored as 32-bit floats without clipping or post-mixing renormalization. Low SNRs can exceed the $[-1, 1]$ range, but clipping is non-linear and would corrupt the intended power ratio.

These are nominal full-band SNRs. Spectral diagnostics were recorded over 25 Hz–8 kHz, but the three categories were not equalized for in-band power, so differences between categories partly reflect their differing spectral distributions and should not be interpreted as differences in perceptual harmfulness alone.

\subsection{Evaluation}
\subsubsection{Primary Metrics}

Each three-second window was one example; the pipeline retained one window per source recording, so no recording received extra weight. Macro-F1 was the primary metric: F1 is computed per instrument and averaged, giving every instrument equal weight regardless of class imbalance. For noise evaluation, macro-F1 was calculated separately for each model, noise category, SNR, and realization.

Clean performance retained was defined as the proportion of clean macro-F1 remaining under a noisy condition:
\begin{equation}
Q_{m,n,s,r}
=
\frac{
F_{m,n,s,r}
}{
F_{m,\mathrm{clean}}
},
\label{eq:retention}
\end{equation}
where \(F_{m,n,s,r}\) is the macro-F1 of model \(m\) under noise category \(n\), SNR \(s\), and realization \(r\), and \(F_{m,\mathrm{clean}}\) is that model's clean macro-F1. A value of \(1\) indicates no degradation; \(0.5\) indicates half the clean macro-F1 retained. The two realization scores were averaged at each noise-category-SNR point, giving \(\overline{Q}_{m,n,s}=(Q_{m,n,s,0}+Q_{m,n,s,1})/2\).

Overall robustness was summarized using the normalized area under the retention-versus-SNR curve. For ascending SNR values $s_1<\cdots<s_J$, normalized AUC was calculated using the trapezoidal rule:
\begin{equation}
A_{m,n}
=
\frac{1}{s_J-s_1}
\sum_{j=1}^{J-1}
\left(s_{j+1}-s_j\right)
\frac{
\overline{Q}_{m,n,s_j}
+
\overline{Q}_{m,n,s_{j+1}}
}{2}.
\label{eq:retention_auc}
\end{equation}
An AUC of $1$ indicates that clean macro-F1 was retained across the complete SNR range. Paired model comparisons were matched by test window, noise category, SNR, and realization, and reported as macro-F1 of model $A$ minus model $B$.


\subsubsection{Secondary Failure Analysis}

CNN and CRNN were compared at 20 dB under white, human non-speech, and environmental noise, with the two realizations analyzed separately. Differences are CNN macro-F1 minus CRNN macro-F1, and positive values favor CNN. Uncertainty was estimated with 2,000 paired bootstrap resamples of complete instrument-pitch groups, drawn with replacement, retaining every test window belonging to each selected group. Both models were evaluated on the same resampled groups. The 95\% percentile-bootstrap interval is the 2.5th and 97.5th percentiles of the resampled differences. Associations between recall loss and five training-derived acoustic summaries were evaluated using Spearman correlations with 100,000 instrument-label permutations. Recall loss was calculated separately for each model, instrument, noise category, SNR, and realization as clean recall minus noisy recall. These associations were exploratory, not confirmatory tests.

To examine whether acoustically similar instruments became confused under noise, each instrument was represented by the mean of its standardized 88-dimensional training-set feature vectors, and acoustic distance between two instruments as the Euclidean distance between these mean vectors. Symmetric confusion for a pair is the mean of the two row-normalized error rates between them. Clean pairwise confusion was subtracted from each noisy value, and the resulting noise-induced confusion was summarized across the complete SNR range using normalized curve area. Spearman correlation related acoustic distance to noise-induced confusion for each model-noise pair. The two realization-specific curve areas were averaged as $\overline{A}_{m,a,b,n}^{\mathrm{conf}}$ for testing, giving the observed relationship
\begin{equation}
\rho_{m,n}
=
\operatorname{Spearman}
\left(
\left\{D_{a,b}\right\},
\left\{\overline{A}_{m,a,b,n}^{\mathrm{conf}}\right\}
\right),
\qquad a<b.
\label{eq:distance_confusion_correlation}
\end{equation}
Statistical significance was estimated using 100,000 permutations in which complete instrument identities were permuted together across matrix rows and columns, preserving dependencies among the 66 instrument pairs. Benjamini-Hochberg(BH) correction was applied across the declared family of 18 comparisons (six models $\times$ three noise categories) at $q<0.05$.


\section{Results}

\subsection{Clean Classification Performance}
Clean macro-F1 spanned 2.0 points across all six models
(Table~\ref{tab:auc}), from 0.9708 (CNN) to 0.9908 (AST), which also
achieved the highest clean accuracy (0.9912). Accuracy and macro-F1 never differed by more than 0.002 for any model. SVM, AST, MERT, and PANNs were evaluated using single fitted runs; CNN and CRNN are five-seed soft-vote ensembles whose single-seed test macro-F1 was $0.962 \pm 0.021$ and
$0.973 \pm 0.003$ respectively.

\subsection{Noise Robustness}
AST had the highest retention AUC under all three noise categories
(Fig.~\ref{fig:curves}, Table~\ref{tab:auc}). At 20 dB under white
noise AST retained 72.3\% of its clean macro-F1 and SVM 10.4\% (61.9-point difference). The lowest AUC was SVM's under white and human
non-speech noise and CNN's under environmental noise.

%%%%%%%%%%%%%%%%%%%%%%%%%%%%%%%%%%%%%%%%%%%%%%%%%%%%%%%%%%%%%%%%%%%%
%% Floats declared HERE, inside Results. A figure* cannot be placed
%% on its declaration page, so the earliest legal slot is the top of
%% the NEXT page -- which is what we want. Do not move these earlier.
%% [!b] + stfloats is deliberate: with [!t] the figure lands above two
%% paragraphs of Methods, before the Results heading exists.
%%%%%%%%%%%%%%%%%%%%%%%%%%%%%%%%%%%%%%%%%%%%%%%%%%%%%%%%%%%%%%%%%%%%

\begin{figure*}[!t]
\centering
\includegraphics[width=\textwidth]{fig6d_retention_compact.png}
\caption{Retention relative to clean macro-F1 versus SNR for all six
models across the three noise categories. Shaded bands span the two
noise realizations.}
\label{fig:curves}
\end{figure*}

\begin{table}[H]
\caption{Clean macro-F1 and retention AUC by noise category. Bold
marks the best value in each column.}
\label{tab:auc}
\centering
\footnotesize
\begin{tabular}{lcccc}
\hline
Model & Clean & White & Human & Environ. \\
\hline
SVM   & 0.9788 & 0.259 & 0.649 & 0.578 \\
CNN   & 0.9708 & 0.357 & 0.653 & 0.505 \\
CRNN  & 0.9738 & 0.330 & 0.655 & 0.515 \\
MERT  & 0.9798 & 0.573 & 0.755 & 0.686 \\
PANNs & 0.9885 & 0.404 & 0.778 & 0.726 \\
AST   & \textbf{0.9908} & \textbf{0.636} & \textbf{0.791} & \textbf{0.753} \\
\hline
\end{tabular}
\end{table}

Model ordering was not preserved across categories, and only AST held its rank under all three. The two noise realizations agreed closely: median retention
range between draws was 0.5 percentage points across all conditions
and 0.6 points at 20 dB, reaching at most 3.6 points, far below the
61.9-point spread between models.

\subsection{Paired Comparisons}
Comparisons were window- and realization-matched. At 20 dB white
noise, CNN exceeded CRNN by 0.1122 and 0.1130 macro-F1; 95\%
pitch-group bootstrap intervals $[0.0560, 0.1625]$ and $[0.0589,
0.1601]$ both excluded zero. Under human non-speech and environmental noise CRNN was better, with intervals excluding zero in one realization but not the other. None of six sign tests survived BH correction.

\begin{figure}[H]
\centering
\includegraphics[width=\columnwidth]
{fig_recall_loss_heat_column.pdf}
\caption{Per-instrument recall-loss AUC across the SNR range, separated
by model and noise category. Instruments are ordered by mean recall
loss across models and categories; all panels share the same colour
scale.}
\label{fig:recall_loss}
\end{figure}

\subsection{Instrument-Specific Failures}
Averaged across models and noise categories, tuba lost the most recall (Fig.2, AUC 0.599) and violin the least (0.215). The pattern was partly model-specific: tuba was the most vulnerable instrument for five of six models under human non-speech noise, but the leading failure class differed under white and environmental noise. The largest single pairwise increase was flute-oboe for the SVM under white noise (0.342). Acoustically closer instrument pairs tended to show greater noise-induced confusion: seven of the 18 predeclared model-by-noise tests were significant after Benjamini-Hochberg correction, most strongly for AST under human non-speech noise ($\rho = -0.607$, permutation $p = 1.0 \times  10^{-5}$, adjusted
$q = 1.8 \times 10^{-4}$).
\begin{figure}[!t]
\centering
\includegraphics[width=\columnwidth]
{fig_distance_confusion_column.pdf}
\caption{Relationship between acoustic distance and noise-induced
confusion. (A) Spearman correlations for all 18 model--noise
comparisons; outlined and starred cells remained significant after BH
correction. (B) Instrument pairs underlying the strongest association,
AST under human non-speech noise.}
\label{fig:distance_confusion}
\end{figure}


\FloatBarrier

\section{Discussion}
% TO BE WRITTEN (~800 words). Results above is deliberately free of % interpretation so that it is not duplicated here.
Results show that clean accuracy is not indicative of robustness under noise. While the models’ macro-F1 and accuracy ratings conformed closely with each other when tested on unperturbed audio, their performances were very different under noise. There were significant variations in model performance ranking for the rest of the models: changes in model ranking were observed across noise types and noise levels. Even still, results generally tended towards AST being the most accurate and robust throughout, while SVM and CNN degraded the quickest. One key takeaway is that clean-trained AST is most noise robust under the conditions evaluated in this study, but these results cannot substantiate conclusions of model performance under noise aware training. There could plausibly be a correlation, but future work is necessary.

\subsection{Limitations}
First, the Philharmonia dataset contains only single-instrument,
single-note recordings. 97.3\% of recordings were shorter than three seconds and were tiled to fill the window, which could confound content with repetition. This issue was explored by evaluating model performance when learning repetition, but only achieved $\sim$20\% accuracy on clean audio. Given this finding, this confound seemed unlikely but its possibility could not be ruled out.  Second, the noise corpus does not cover recording or studio effects such as reverberation. Third, only two noise realizations were generated, giving a sensitivity check but not a reliable variance estimate. Fourth, nominal SNR is full-band, and the categories concentrate energy differently across frequency, so they may mask instruments unequally, making reliable, direct comparison of noise loudness unattainable. Lastly, models were trained on clean data only, so this study measures robustness without any noise adaptation.

\section{Conclusion}
Our research compares and ranks six musical-instrument classification systems under controlled white, human, and environmental noise degradation. Despite similar performance on clean recordings, AST retained the most clean macro-F1 across the tested conditions, while SVM degraded earliest under white noise and CNN under environmental. Failures were also uneven across instruments: tuba, oboe, and trumpet showed the most recall loss. Acoustically similar instrument pairs became more confused under noise in 7 of 18 model-noise tests after correction, most consistently under recorded noise. These findings demonstrate that clean accuracy alone is insufficient for selecting an instrument classifier for noisy conditions. Robustness must also be evaluated across multiple noise categories and severities. Future work should extend this benchmark to multiple datasets, polyphonic music, independently recorded environmental noise, additional noise realizations, and noise-aware training.

\section*{Acknowledgments}

This work was created in affiliation with Boston University's RISE program. We are grateful to Boston University for this opportunity as well as our instructors Dr. Eugene Pinsky and Dr. Indrajit Kalita for their dedication and hard work.

\begin{thebibliography}{00}
\bibitem{b1} S. McAdams, "Perspectives on the contribution of timbre to musical structure," Computer Music Journal, vol. 23, no. 3, pp. 85–102, 1999.
\bibitem{b2} M. Blaszke and B. Kostek, "Musical instrument identification using deep learning approach," Sensors, vol. 22, no. 8, art. 3033, 2022.
\bibitem{b3} J. Marques and P. J. Moreno, "A study of musical instrument classification using Gaussian mixture models and support vector machines," Cambridge Research Laboratory, Compaq Computer Corp., Tech. Rep. CRL 99/4, 1999.
\bibitem{b4} P. Herrera-Boyer, G. Peeters, and S. Dubnov, "Automatic classification of musical instrument sounds," Journal of New Music Research, vol. 32, no. 1, pp. 3–21, 2003.
\bibitem{b5} A. Livshin and X. Rodet, "Instrument recognition beyond separate notes — indexing continuous recordings," in Proc. Int. Computer Music Conf. (ICMC), Miami, FL, USA, 2004.
\bibitem{b6} T. Park and T. Lee, "Musical instrument sound classification with deep convolutional neural network using feature fusion approach," arXiv:1512.07370, 2015.
\bibitem{b7} S. Prabavathy, V. Rathikarani, and P. Dhanalakshmi, "Classification of musical instruments sound using pre-trained model with machine learning techniques," Asian Journal of Electrical Sciences, vol. 9, no. 1, pp. 45–48, 2020.
\bibitem{b8} J. Zheng, M. Cao, and C. Zhang, "ICKAN: A deep musical instrument classification model incorporating Kolmogorov–Arnold network," Scientific Reports, vol. 15, art. 21573, 2025.
\bibitem{b9} L. Zhong, E. Cooper, J. Yamagishi, and N. Minematsu, "Exploring isolated musical notes as pre-training data for predominant instrument recognition in polyphonic music," in Proc. APSIPA ASC, 2023.
\bibitem{b10} V. Deng, C. Wang, G. Richard, and B. McFee, "Investigating the sensitivity of pre-trained audio embeddings to common effects," in Proc. IEEE ICASSP, Hyderabad, India, Apr. 2025.
\bibitem{b11} S. Gonzales, D. Garcia, A. Sanchez, H. San Blas, and G. Gonzales, "Evaluating pre-trained audio embeddings for the classification of traditional Portuguese musical instruments," Machine Learning with Applications, vol. 24, art. 100910, 2026.
\bibitem{b12} Philharmonia, “Sound samples,” Philharmonia. [Online]. Available: https://philharmonia.co.uk/resources/sound-samples/. [Accessed: July 25, 2026].
\bibitem{b13} K. J. Piczak, "ESC: Dataset for environmental sound classification," in Proc. 23rd ACM Int. Conf. Multimedia, 2015, pp. 1015–1018.
\bibitem{b14} J. Thiemann, N. Ito, and E. Vincent, "The Diverse Environments Multi-channel Acoustic Noise Database (DEMAND)," in Proc. Int. Congr. Acoustics (ICA), Montreal, Canada, 2013.
\bibitem{b15}Y. Li et al., “MERT: Acoustic music understanding model with large-scale self-supervised training,” arXiv preprint arXiv:2306.00107, 2023, doi: 10.48550/arXiv.2306.00107
\bibitem{b16}Y. Gong, Y.-A. Chung, and J. Glass, “AST: Audio Spectrogram Transformer,” Proc. Interspeech, 2021, doi: 10.48550/arXiv.2104.01778.
\bibitem{b17} Q. Kong, Y. Cao, T. Iqbal, Y. Wang, W. Wang, and M. D. Plumbley, “PANNs: Large-scale pretrained audio neural networks for audio pattern recognition,” arXiv preprint arXiv:1912.10211, 2020, doi: 10.48550/arXiv.1912.10211

\end{thebibliography}
\end{document}
